%! TEX program = lualatex

%%%%%%%%%%%%%%%%%%%%%%%%%%%%%%%%%%%%%%%%%
% Professional Formal Letter
% LaTeX Template
% Version 2.0 (12/2/17)
%
% This template originates from:
% http://www.LaTeXTemplates.com
%
% Authors:
% Brian Moses
% Vel (vel@LaTeXTemplates.com)
%
% License:
% CC BY-NC-SA 3.0 (http://creativecommons.org/licenses/by-nc-sa/3.0/)
%
%%%%%%%%%%%%%%%%%%%%%%%%%%%%%%%%%%%%%%%%%

%-------------------------------------------------------------------------------
%	PACKAGES AND OTHER DOCUMENT CONFIGURATIONS
%-------------------------------------------------------------------------------

\documentclass[11pt, a4paper]{letter} % Set the font size (10pt, 11pt and 12pt) and paper size (letterpaper, a4paper, etc)

\input{structure.tex} % Include the file that specifies the document structure

\longindentation=0pt % Un-commenting this line will push the closing "Sincerely," and date to the left of the page

%-------------------------------------------------------------------------------
%	YOUR INFORMATION
%-------------------------------------------------------------------------------

\Who{Simon A. Babayan} % Your name

\Title{, PhD} % Your title, leave blank for no title

\authordetails{
	School of Biodiversity, One Health,\\ and Veterinary Medicine\\ % Your department/institution
	Graham Kerr Building\\ % Your address
	Glasgow, G12 8QQ\\ % Your city, zip code, country, etc
	Email: simon.babayan@glasgow.ac.uk\\ % Your email address
}

%-------------------------------------------------------------------------------
%	HEADER CONTENTS
%-------------------------------------------------------------------------------

\logo{UoG_logo.png} % Logo filename, your logo should have square dimensions (i.e. roughly the same width and height), if it does not, you will need to adjust spacing within the HEADER STRUCTURE block in structure.tex (read the comments carefully!)

\headerlineone{University} % Top header line, leave blank if you only want the bottom line

\headerlinetwo{\emph{of} Glasgow} % Bottom header line

%-------------------------------------------------------------------------------

\begin{document}

%-------------------------------------------------------------------------------
%	TO ADDRESS
%-------------------------------------------------------------------------------

\begin{letter}{
    Editorial Board of \emph{PLoS Biology}.
}

%-------------------------------------------------------------------------------
%	LETTER CONTENT
%-------------------------------------------------------------------------------

\opening{Dear Editor,}

We are pleased to submit our final revision of the manuscript entitled \textbf{``Environmental drivers of low vaccine responsiveness in a lab-to-wild rodent model''} for publication in PLoS Biology.

Our manuscript addresses a critical challenge in vaccinology: the significant gap between vaccine performance in controlled laboratory conditions versus diverse real-world environments. This is most manifest in the high failure rate of vaccines late into their development cycle, e.g. in human efficacy trials (Phase III). Even successful vaccines exhibit substantial reductions in efficacy when deployed in populations that differ from those in which they were developed, particularly in rural versus urban settings. Despite widespread recognition of this problem, the specific environmental factors driving this vaccine hyporesponsiveness remain poorly understood and difficult to quantify.

In this study, we present a novel approach to disentangle the complex interplay of habitat, diet, sex, and parasite burden on vaccine-specific antibody responses. Using paired cohorts of laboratory-reared and wild wood mice (\textit{Apodemus sylvaticus}), we demonstrate that circulating vaccine-specific IgG1 antibodies were 47\% lower in wild populations compared to laboratory-reared conspecifics. Contrary to expectations, our results reveal that dietary supplementation—often considered beneficial—actually reduced vaccine responsiveness, which raises questions about the potentially heightened responsiveness to immunisation under resource scarcity. Additionally, we show that helminth infection negatively affects vaccine-specific antibody production once habitat, diet, and sex are accounted for.

Our findings are particularly significant because they:

\begin{enumerate}
\item Provide quantitative evidence that laboratory settings may systematically overestimate vaccine performance by failing to capture natural environmental variability
\item Establish a robust, causally-explicit modelling approach through structural causal models (SCMs) to identify and quantify factors driving differential vaccine responsiveness
\item Suggest that current efficacy thresholds for advancing vaccines through clinical development stages may need reconsideration
\item Offer a translational model system for testing interventions to improve vaccine effectiveness in different environmental contexts
\end{enumerate}

This work aligns with PLOS Biology's mission of publishing significant advances across biological sciences by bridging the fields of ecology, immunology, and public health. The interdisciplinary approaches we employ address questions of fundamental biological importance while offering insights that could substantially improve vaccination strategies in diverse human populations.

All authors have approved the manuscript and agree with its submission to PLOS Biology. The manuscript has not been published elsewhere and is not under consideration by another journal. We have no conflicts of interest to declare.

We suggest the following experts as potential reviewers for our manuscript:

\begin{enumerate}
\item Reviewer name, institution, email
\item Reviewer name, institution, email
\item Reviewer name, institution, email
\item Reviewer name, institution, email
\end{enumerate}

Thank you for considering our manuscript. We look forward to your response.

\closing{Yours sincerely, on behalf of my co-authors,}

%-------------------------------------------------------------------------------
%	OPTIONAL RIGHT-HAND FOOTNOTE
%-------------------------------------------------------------------------------

% Uncomment the 4 lines below to print a footnote with custom text
\def\thefootnote{}
\def\footnoterule{\hrule}
\def\thefootnote{\arabic{footnote}}

%-------------------------------------------------------------------------------

\end{letter}

\end{document}
